% Methods for the J-BHI paper. Hand-written prose; every result value is a macro from
% generated/c2c1_numbers.tex or generated/c2c2_numbers.tex. Constants typed here are protocol
% parameters fixed in the registrations (osf.io/s6w9c, osf.io/4pzxa), not results.

\section{\textbf{Materials and Methods}}\label{sec:methods}

\subsection{\textbf{Problem Formulation}}\label{sec:formulation}
Let $i = 1,\dots,n$ index Parkinson's disease (PD) donors, each with an induced pluripotent stem
cell (iPSC) line differentiated towards dopaminergic (DA) neurons, and let $j = 1,\dots,n_i$ index
the donor's clinical visits at times $t_{ij}$ (years since diagnosis). From the day-65 culture of
donor $i$ we compute a scalar cellular phenotype $P_i$, and from the clinic we observe a score
$Y_{ij}$. The question is whether $P_i$ predicts the rate of change of $Y$, i.e.\ whether the
interaction coefficient $\beta_3$ in a longitudinal model of $Y_{ij}$ (Section~\ref{sec:lmm}) is
non-zero.

A cellular phenotype can fail for reasons unrelated to biology: the processing may not reproduce,
the phenotype may not be stable across replicate differentiations, the labelled cells may not be
the intended cell type, or the clinical data may be too sparse to estimate slopes. We therefore
formulate the study as a sequence of \emph{gates}. Each gate $G_k$ is a statistic $T_k$ computed
without clinical outcomes and a threshold $\tau_k$ fixed in advance,
\begin{gather}
G_k(P) = \mathbb{1}\!\left[T_k(P) \ge \tau_k\right],
\label{eq:gates}\\
\mathcal{P}_{\mathrm{link}} = \Big\{P \in \mathcal{P} : \textstyle\prod_k G_k(P) = 1\Big\},
\label{eq:linkset}
\end{gather}
and only phenotypes in $\mathcal{P}_{\mathrm{link}}$ are linked to clinical data. Study-level
gates (pipeline reproduction, clinical sufficiency) stop the analysis if they fail; phenotype-level
gates (stability, coherence) remove a phenotype and report it as ``not tested''. All thresholds,
models and decision rules were registered on OSF before linkage (registration~1:
osf.io/s6w9c; registration~II: osf.io/4pzxa), and each analysis step ran from a git tag frozen
before it touched the data \cite{nosek2018prereg}. Fig.~\ref{fig:gates} shows the gate sequence.

\subsection{\textbf{Data}}\label{sec:data}
\emph{Cellular data:} FOUNDIN-PD \cite{bressan2023foundin} provides day-65 single-cell RNA
sequencing (scRNA-seq) of \cconeRuns{} runs from \cconeLines{} iPSC lines: one run for each of
79 donor lines and four replicate runs of one control line, which serve as technical
replicates. Lines were differentiated in five batches; the last two used a different
medium (E6). We also use day-0 and day-65 bulk RNA-seq and NeuroChip genotypes (after quality
control: call rate $\ge 0.95$, MAF $\ge 0.05$, Hardy--Weinberg $p \ge 10^{-6}$; \cctwoArraySNPs{}
autosomal SNPs).

\emph{Clinical data.} PPMI \cite{marek2011ppmi} provides longitudinal follow-up of the same donors.
The primary outcome is the OFF-state MDS-UPDRS Part~III total \cite{goetz2008mdsupdrs}; a visit is
OFF if the donor was untreated or the exam was recorded in the OFF state, and repeated exams on one
date are averaged. Secondary outcomes are the mean left/right putamen specific binding ratio (SBR)
from DaTscan, CSF $\alpha$-synuclein seed amplification (SAA) status (earliest valid assay)
\cite{siderowf2023saa}, and the neuronal $\alpha$-synuclein disease integrated stage (NSD-ISS;
latest record, integer stage) \cite{simuni2024nsdiss}. Only aggregate results leave the
analysis machine.

\subsection{Single-Cell Processing and Cell Labels}\label{sec:cells}
\emph{Quality control.} For cell $c$ with UMI counts $u_{gc}$ over genes $g$, let $n_c$ be the
number of detected genes, $U_c = \sum_g u_{gc}$ and $m_c$ the mitochondrial UMI share. A cell is
retained if
\begin{equation}
500 \le n_c \le 8000, \qquad U_c \ge 1000, \qquad m_c \le 20\%.
\label{eq:qc}
\end{equation}

\emph{Module scores.} With normalised expression $x_{gc} = \ln\!\big(1 + 10^4 u_{gc}/U_c\big)$, the
score of cell $c$ for marker set $M$ is the mean marker expression minus that of a
control set $R_M$ drawn to match the markers' expression distribution (50 genes from each of 25
expression bins, fixed seed):
\begin{equation}
s_c(M) = \frac{1}{|M|}\sum_{g\in M} x_{gc} \;-\; \frac{1}{|R_M|}\sum_{g\in R_M} x_{gc}.
\label{eq:score}
\end{equation}
Each cell is assigned to the compartment $k^\star(c) = \arg\max_k s_c(M_k)$ among neuron,
progenitor, proliferating and ependymal-like marker sets, or left unassigned if all scores are
negative.

\emph{DA-neuron labels.} Registration~1 used the rule ``neuron with $u_{\mathrm{TH},c} \ge 1$''
(old label). Registration~II uses a nine-marker DA module
$M_{\mathrm{DA}} = \{$TH, DDC, SLC6A3, KCNJ6, SLC18A2, NR4A2, LMX1B, FOXA2, EN1$\}$:
\begin{equation}
\mathrm{DA}(c) = \mathbb{1}\!\left[k^\star(c) = \text{neuron}\right]\cdot
\mathbb{1}\!\left[s_c(M_{\mathrm{DA}}) > t^\star\right].
\label{eq:dalabel}
\end{equation}
The single threshold is calibrated on public data only, as the value that best reproduces the
published DA-neuron fractions $f^{\mathrm{pub}}_\ell$ of the donor lines $\ell$:
\begin{equation}
t^\star = \arg\max_{t \in \mathcal{T}}\;
\operatorname{corr}_\ell\!\left(\operatorname{logit} f_\ell(t),\, \operatorname{logit} f^{\mathrm{pub}}_\ell\right),
\label{eq:calib}
\end{equation}
over the grid $\mathcal{T} = \{-0.10, -0.075, \dots, 0.60\}$, where $f_\ell(t)$ is the share of QC cells of line $\ell$ labelled DA at threshold $t$, and
fractions are clipped to $[0.005, 0.995]$ before the logit.

\subsection{Phenotype Construction}\label{sec:phenotypes}
For a cell set $A$ (all QC cells, or DA neurons) of line $\ell$, the pseudobulk count of gene $g$
and its log expression are
\begin{equation}
C_{g\ell} = \sum_{c \in A_\ell} u_{gc}, \qquad
y_{g\ell} = \log_2\!\Big(1 + 10^6\, C_{g\ell} \big/ \textstyle\sum_{g'} C_{g'\ell}\Big).
\label{eq:pseudobulk}
\end{equation}
A single-gene phenotype is $y_{g\ell}$ itself (SNCA). A pathway phenotype averages gene-wise
standardised expression over the set $G'$ of pathway genes with mean CPM $\ge 10$:
\begin{equation}
S_\ell(G) = \frac{1}{|G'|}\sum_{g \in G'} \frac{y_{g\ell} - \bar{y}_g}{\sigma_g},
\label{eq:geneset}
\end{equation}
with $\bar{y}_g, \sigma_g$ the mean and population SD over the 79 donor lines (HALLMARK oxidative
phosphorylation, OXPHOS; KEGG lysosome).

\emph{Registration~1} uses all-cell phenotypes: P1, the logit of the published DA-neuron fraction;
P2, SNCA; P3, OXPHOS; P4, lysosome. Because all-cell expression mixes cell types, P2--P4 are
residualised on composition,
\begin{gather}
\mathbf{P} = \big(\mathbf{I} - \mathbf{W}(\mathbf{W}^\top\mathbf{W})^{-1}\mathbf{W}^\top\big)\,\tilde{\mathbf{p}},
\label{eq:resid}\\
\mathbf{W} = \big[\mathbf{1},\ \operatorname{logit}\mathbf{f}_{\mathrm{DA}},\ \operatorname{logit}\mathbf{f}_{\mathrm{prog}}\big],
\nonumber
\end{gather}
where $\tilde{\mathbf{p}}$ is the raw phenotype over lines and $\mathbf{f}$ are the published
DA-neuron and progenitor fractions. \emph{Registration~II} computes the phenotypes inside DA neurons
(\ref{eq:dalabel}): D1, the logit DA fraction; D2, SNCA; D3, OXPHOS; D4, lysosome. These need no
composition adjustment; the residualised versions are a sensitivity analysis. Every phenotype is
standardised to mean 0 and SD 1 over the donors of each model, so effects are per SD.

\subsection{Gates Before Linkage}\label{sec:prelink}
\emph{Reproduction (A0-6).} Our per-line fraction of TH$^+$ cells must correlate with the published
value at Pearson $r \ge 0.80$.

\emph{Positive control (A0-4).} A known biological signal must be recoverable: we replicate the
published day-0 classifier of differentiation efficiency \cite{bressan2023foundin}, an elastic net
on day-0 bulk expression $\mathbf{x}_\ell$ predicting high vs.\ low DA yield $y_\ell \in \{0,1\}$,
\begin{equation}
\min_{b_0,\mathbf{b}}\; -\frac{1}{n}\sum_\ell \big[y_\ell \eta_\ell - \ln(1 + e^{\eta_\ell})\big]
+ \lambda\Big(\tfrac{1-\alpha}{2}\|\mathbf{b}\|_2^2 + \alpha\|\mathbf{b}\|_1\Big),
\label{eq:enet}
\end{equation}
with $\eta_\ell = b_0 + \mathbf{x}_\ell^\top\mathbf{b}$, $\alpha = 0.5$, $\lambda = 0.22$. It passes
if the mean AUC over 100 repeats of 5-fold cross-validation is $\ge 0.75$ with permutation
$p < 0.05$ (200 label permutations), and is ``partial'' at AUC $\ge 0.65$.

\emph{Technical stability (A0-2, G-icc).} The control line's replicate runs estimate within-line
noise. With $S_d^2$ the variance of the phenotype over the $n_d = 79$ donor lines and $S_w^2$ the
variance over the $n_w = 4$ control runs, the intraclass correlation and its 95\% CI are
\begin{gather}
\widehat{\mathrm{ICC}} = 1 - \frac{S_w^2}{S_d^2},
\label{eq:icc}\\
\mathrm{CI}_{95} = \Big[\,1 - F_{0.975}\,\tfrac{S_w^2}{S_d^2},\;\; 1 - F_{0.025}\,\tfrac{S_w^2}{S_d^2}\,\Big],
\label{eq:icc_ci}
\end{gather}
where $F_q$ is the $q$ quantile of $F(n_d - 1,\, n_w - 1)$, since
$(1-\mathrm{ICC})\,S_d^2/S_w^2 \sim F(n_d-1, n_w-1)$. A phenotype passes at
$\widehat{\mathrm{ICC}} \ge 0.5$ \cite{koo2016guideline}. Because the control runs may share
batch effects, registration~II also requires an across-batch ICC $\ge 0.40$ for the same phenotype
in an independent dataset of 199 lines differentiated in many batches \cite{jerber2021population}.

\emph{Cell sufficiency (A0-3, G-cells).} A run needs $\ge 1{,}000$ QC cells (registration~1) and
$\ge 100$ DA neurons (registration~II).

\emph{Clinical sufficiency (A0-1).} Slopes need repeated visits:
\begin{equation}
N_{\mathrm{link}} = \sum_i \mathbb{1}\!\left[n_i^{\mathrm{OFF}} \ge 3\ \wedge\
\max_j t_{ij} - \min_j t_{ij} \ge 2\right] \ge 30.
\label{eq:a01}
\end{equation}

\subsection{Genetic Positive Control (G-eQTL)}\label{sec:geqtl}
If the new DA-neuron pseudobulks reflect DA neurons, they should carry the genetic regulation
known from DA neurons. Lead eQTLs (gene $g$, SNP $v_g$, effect $b^{\mathrm{J}}_g$) were taken from
day-52 untreated DA neurons of Jerber et al.\ \cite{jerber2021population}, one lead per gene,
Benjamini--Hochberg (BH) FDR $< 5\%$ on the gene-level permutation $p$. In the 79 FOUNDIN donor
lines, each pair is tested by
\begin{equation}
y_{g\ell} = a_g + b_g\, d_{v_g\ell} + \sum_{k=1}^{3}\gamma_k\,\mathrm{GPC}_{k\ell}
+ \sum_{k=1}^{5}\delta_k\,\mathrm{EPC}_{k\ell} + \varepsilon_{g\ell},
\label{eq:eqtl}
\end{equation}
where $d \in [0,2]$ is the allele dosage aligned to Jerber's assessed allele, GPC are genotype
principal components and EPC are principal components of the 2,000 most variable genes. SNPs are
matched on GRCh37 position and alleles, and strand-ambiguous SNPs are dropped. Over the $m$ testable
pairs we compute the sign concordance and Storey's $\pi_1$ \cite{storey2003qvalue}:
\begin{gather}
\kappa = \frac{1}{m}\sum_{g=1}^{m}\mathbb{1}\!\left[\operatorname{sign}\hat{b}_g = \operatorname{sign} b^{\mathrm{J}}_g\right],
\label{eq:kappa}\\
\hat{\pi}_1 = 1 - \frac{\#\{g : p_g > \lambda\}}{m(1-\lambda)},
\label{eq:pi1}
\end{gather}
with $\lambda = 0.5$. A null distribution comes from $B = 200$ sets that pair the same genes with
random SNPs on other chromosomes in the same MAF bin. The gate passes if
$\kappa \ge 0.70$, $\hat\pi_1 \ge 0.30$, and both exceed the 95th percentile of their null.

\subsection{Longitudinal Linkage Model}\label{sec:lmm}
For donor $i$ at visit $j$, with standardised phenotype $P_i$ and covariate vector $\mathbf{c}_i$
(age at onset, sex, genetic PD, polygenic risk score, medium E6, disease duration at PPMI baseline),
\begin{multline}
Y_{ij} = \beta_0 + \beta_1 t_{ij} + \beta_2 P_i + \beta_3 P_i t_{ij}
+ \boldsymbol{\gamma}^\top\mathbf{c}_i + \boldsymbol{\delta}^\top\mathbf{c}_i\, t_{ij}\\
+ b_{0i} + b_{1i} t_{ij} + \varepsilon_{ij},
\label{eq:lmm}
\end{multline}
with $(b_{0i}, b_{1i})^\top \sim \mathcal{N}(\mathbf{0}, \boldsymbol{\Sigma}_b)$ and
$\varepsilon_{ij} \sim \mathcal{N}(0, \sigma^2)$. The parameter of interest is $\beta_3$: the change
in progression rate (UPDRS~III points per year) per SD of the phenotype. In matrix form,
$\mathbf{y}_i \sim \mathcal{N}(\mathbf{X}_i\boldsymbol{\beta}, \mathbf{V}_i)$ with
$\mathbf{V}_i = \mathbf{Z}_i\boldsymbol{\Sigma}_b\mathbf{Z}_i^\top + \sigma^2\mathbf{I}$ and
$\mathbf{Z}_i = [\mathbf{1}, \mathbf{t}_i]$. The variance parameters
$\boldsymbol{\psi} = (\boldsymbol{\theta}, \sigma)$, with
$\boldsymbol{\Sigma}_b = \sigma^2\boldsymbol{\Lambda}_\theta\boldsymbol{\Lambda}_\theta^\top$ and
$\boldsymbol{\Lambda}_\theta$ lower triangular, maximise the restricted log-likelihood
\begin{multline}
\ell_R(\boldsymbol{\psi}) = -\tfrac{1}{2}\Big[\sum_i \ln|\mathbf{V}_i|
+ \ln\Big|\sum_i \mathbf{X}_i^\top\mathbf{V}_i^{-1}\mathbf{X}_i\Big|\\
+ \sum_i (\mathbf{y}_i - \mathbf{X}_i\hat{\boldsymbol\beta})^\top\mathbf{V}_i^{-1}(\mathbf{y}_i - \mathbf{X}_i\hat{\boldsymbol\beta})\Big],
\label{eq:reml}
\end{multline}
where $\hat{\boldsymbol\beta}$ is the generalised least-squares estimate with covariance
$\boldsymbol{\Phi} = (\sum_i \mathbf{X}_i^\top\mathbf{V}_i^{-1}\mathbf{X}_i)^{-1}$. For the contrast
$\mathbf{L}$ that selects $\beta_3$,
\begin{equation}
t = \frac{\mathbf{L}^\top\hat{\boldsymbol\beta}}{\sqrt{\mathbf{L}^\top\hat{\boldsymbol\Phi}\mathbf{L}}},
\qquad
\nu = \frac{2\big(\mathbf{L}^\top\hat{\boldsymbol\Phi}\mathbf{L}\big)^2}
{\mathbf{g}^\top\mathbf{A}\,\mathbf{g}},
\label{eq:satt}
\end{equation}
where $\nu$ is the Satterthwaite degrees of freedom \cite{satterthwaite1946},
$\mathbf{g} = \partial(\mathbf{L}^\top\boldsymbol\Phi\mathbf{L})/\partial\boldsymbol\psi$ and
$\mathbf{A} = 2\mathbf{H}^{-1}$ is the asymptotic covariance of $\hat{\boldsymbol\psi}$ from the
Hessian $\mathbf{H}$ of the REML deviance $-2\ell_R$; both are computed by central finite
differences. If the fit fails to converge, is singular (a random-effect SD below $10^{-4}\sigma$
or $|\mathrm{corr}| > 0.999$) or has a non-positive-definite $\mathbf{H}$, the registered fallbacks
are applied in order: uncorrelated random intercept and slope, then random intercept only. The
same model with $Y$ = putamen SBR (all donors with $\ge 1$ scan) gives H2.

\emph{Cross-sectional outcomes (H3).} SAA status and NSD-ISS stage $Z_i$ are tested by ordinary
least squares, $P_i = \alpha_0 + \alpha_1 Z_i + \boldsymbol\gamma^\top\mathbf{c}_i + e_i$,
with $\nu = n - p$ residual degrees of freedom.

\subsection{Effect Size, Equivalence and Decision Rule}\label{sec:decision}
Each test is reported as a partial correlation,
\begin{equation}
r = \frac{t}{\sqrt{t^2 + \nu}},
\label{eq:partialr}
\end{equation}
which is monotone in $t$. The 90\% CI of the coefficient,
$\hat\beta \pm t_{0.95,\nu}\,\widehat{\mathrm{SE}}$, therefore maps to a 90\% CI
$[r_{\mathrm{lo}}, r_{\mathrm{hi}}]$ of $r$. The test is \emph{equivalent} to zero within the
bound $\Delta = 0.40$ if
\begin{equation}
E = \mathbb{1}\!\left[-\Delta < r_{\mathrm{lo}}\ \wedge\ r_{\mathrm{hi}} < \Delta\right],
\label{eq:tost}
\end{equation}
which is the two one-sided tests (TOST) procedure at $\alpha = 0.05$ \cite{lakens2017equivalence}.
Within a family of $m$ tests, $p$-values are converted to BH $q$-values \cite{benjamini1995fdr},
\begin{equation}
q_{(k)} = \min_{j \ge k}\ \min\!\Big(1,\ \frac{m\,p_{(j)}}{j}\Big),
\label{eq:bh}
\end{equation}
and every test receives one of three classes (for single tests, $q$ is replaced by $p$):
\begin{equation}
\mathrm{class} =
\begin{cases}
\text{supported}, & q < 0.05,\\
\text{equivalent}, & q \ge 0.05,\ E = 1,\\
\text{inconclusive}, & q \ge 0.05,\ E = 0.
\end{cases}
\label{eq:class}
\end{equation}

\emph{Coherence (A0-5).} A pathway score is tested only if its member genes agree with it. Refitting
(\ref{eq:lmm}) with each member gene $g \in G'$ in place of $P$,
\begin{equation}
\pi_{\mathrm{coh}} = \frac{1}{|G'|}\sum_{g \in G'}
\mathbb{1}\!\left[\operatorname{sign}\hat\beta_3^{(g)} = \operatorname{sign}\hat\beta_3^{(P)}\right] \ge 0.80.
\label{eq:a05}
\end{equation}

\subsection{Hypotheses}\label{sec:hypotheses}
\emph{Registration~1} tests the all-cell phenotypes retained by A0-2 against four outcomes, each a
BH family: H1, UPDRS~III slope; H2, SBR slope; H3, SAA status and NSD-ISS stage.
\emph{Registration~II} has one primary test, H1-DA: $\beta_3$ for D2 on the OFF-state UPDRS~III,
two-sided $\alpha = 0.05$. A secondary BH family holds D4 on UPDRS~III (if A0-5 passes), D2 and D4 on
SBR slope, SAA and NSD-ISS. A separate single test, H4, asks whether the registration-1 exploratory
OXPHOS signal survives a change of assay: HALLMARK OXPHOS in day-65 bulk RNA, composition-adjusted
as in (\ref{eq:resid}), on UPDRS~III. The registered sensitivity analyses of H1-DA are: old DA label;
composition-adjusted D2; idiopathic PD only; five-level batch dummies in place of the medium
covariate; leave-one-donor-out range of $\hat\beta_3$; and visits within 5 years of diagnosis.

\subsection{Exploratory Batch Decomposition}\label{sec:batchmethod}
After the primary result was seen to weaken under batch adjustment, one post-hoc analysis was
specified and tagged before it ran. With per-donor OLS slopes $\hat s_i$ of UPDRS~III on time
($\ge 3$ OFF visits), batch $b(i) \in \{1,\dots,5\}$ and $x_i$ = standardised D2, it measures
batch structure, $R^2_x = \sum_b n_b(\bar x_b - \bar x)^2 / \sum_i (x_i - \bar x)^2$ (and likewise
for $\hat s$), and splits the association into within- and between-batch parts:
\begin{equation}
r_{W} = \operatorname{corr}\!\left(x_i - \bar x_{b(i)},\ \hat s_i - \bar{\hat s}_{b(i)}\right),
\qquad
r_{B} = \operatorname{corr}_b\!\left(\bar x_b,\ \bar{\hat s}_b\right),
\label{eq:within}
\end{equation}
with $\nu_W = n - 2 - (B - 1)$ for $r_W$. The reading rules were fixed in advance: confounding is
plausible if both batch effects have $p < 0.10$ and $r_B$ has the sign of the overall $r$;
otherwise the association lies mainly within batches if $|r_W| \ge 0.75\,|r|$ with the same sign.

\subsection{Implementation}\label{sec:impl}
All steps are implemented in Python with NumPy, SciPy, pandas and h5py; the REML fit,
Satterthwaite degrees of freedom, elastic net and gene-set scoring were implemented from their
definitions and checked against planted-effect synthetic data. Each step ran from a git tag
created before it saw the relevant data, and failed versions are reported (Section~\ref{sec:deviations}).
